% source: erdos-872/current_state.md (The 5/24 first-hit cover)
% source: erdos-872/lean/README.md

\section{The Exact \texorpdfstring{$5/24$}{5/24} Cover Theorem}\label{sec:524-cover}

This section concerns the unweighted first-hit cover problem on the lower half.
It is logically independent of the game values $\L(n)$, but it isolates the
cleanest exact extremal structure in the repo and supplies a useful geometric
model for later arguments.

\begin{definition}
Let $\taucov(n)$ be the minimum size of a set $H \subseteq \U$ with the
property that every $x \in L$ divides some $u \in H$.
\end{definition}

The extremal cover and packing are explicit.

\begin{theorem}\label{thm:524-cover}
Let
\[
  H_n
  :=
  \{u \in \U : u \equiv 2 \pmod 4\}
  \cup
  \{u \in \U : u > 2n/3,\ u \equiv 0 \pmod 4\},
\]
and
\[
  P_n
  :=
  \big((n/3,n/2]\cap\Z\big)
  \cup
  \{x \in (n/4,n/3]\cap\Z : x \text{ odd}\}.
\]
Then:
\begin{enumerate}
  \item every $x \in L$ divides some $u \in H_n$;
  \item every $u \in \U$ is divisible by at most one element of $P_n$;
  \item $|H_n| = |P_n| = \dfrac{5}{24}n + O(1)$.
\end{enumerate}
Consequently
\[
  \taucov(n) = \frac{5}{24}n + O(1).
\]
\end{theorem}

\begin{proof}[Proof sketch]
For the cover, one checks the lower half by dyadic strips.  If
$x \in (n/3,n/2]$, then $2x \in \U$ and $2x \equiv 2 \pmod 4$, so $2x \in H_n$.
If $x \in (n/4,n/3]$, then either $2x \in \U$ is congruent to $2 \pmod 4$, or
$3x \in \U$ and falls into the top $4$-divisible strip.  The remaining lower
strips are handled similarly: for each $x \in (n/(j+1),n/j]$ with $j\ge 4$, one
of the multipliers $2$ or $3$ moves $x$ into the prescribed portion of $\U$.

For the packing, the interval $(n/3,n/2]$ is easy: if two distinct elements
$x<y$ from that interval both divide $u \in \U$, then $2x>n/2$ and $2y>n/2$,
forcing $u$ to be larger than $n$.  For odd $x \in (n/4,n/3]$, any multiple of
$x$ in $\U$ must be either $2x$ or $3x$, and the key lemma is that
\[
  \frac{u}{x} \in \{2,3\}.
\]
This prevents one upper-half integer from hitting two distinct packed points.

Finally,
\[
\begin{aligned}
  |H_n|
  &=
  \#\{u \in \U : u \equiv 2 \!\!\!\pmod 4\} \\
  &\quad +
  \#\{u \in \U : u > 2n/3,\ u \equiv 0 \!\!\!\pmod 4\} \\
  &=
  \frac{5}{24}n + O(1),
\end{aligned}
\]
and the same floor-sum count gives $|P_n| = 5n/24 + O(1)$.  Since $H_n$ is a
cover and $P_n$ is a packing, optimality follows.
\end{proof}

\begin{remark}
The formal verification covers the structural part of the theorem: the cover
property of $H_n$, the packing property of $P_n$, and the exact inequality
showing that one upper-half integer cannot cover two packed lower-half points.
The explicit count $5n/24+O(1)$ is an elementary floor-sum calculation and is
not the interesting formal bottleneck.
\end{remark}
